\section{Preventing the Misuse of AI for Authoritarian Purposes}
\label{sec:6.4}
An AI system does not need a broken moral sense to become an instrument of control. It can work exactly as designed and still do damage, if the government directing it wants to track its citizens, silence dissent, or manufacture consent it never earned. Facial recognition built to unlock a phone recognizes a protester's face just as well. A recommendation system built to hold attention can be tuned to spread a regime's preferred narrative instead of a user's preferred content. The technology does not change; only the hand on the controls does.

Meeting that risk takes four things in sequence: knowing what authoritarian misuse of AI actually looks like, rather than treating it as hypothetical; building systems technically resistant to that misuse from the start; the policy work that has to do the rest, because no engineering choice by itself can stop a government that owns the infrastructure from disabling whatever safeguard stands in its way; and then the case where the government does not have to disable anything, because it can compel the people who hold the system instead. The last of the four is where the other three run out.
